\documentclass[aps,prd,reprint,amsmath,amssymb,nofootinbib,longbibliography]{revtex4-2}

\usepackage{booktabs}
\usepackage{bm}

\makeatletter
\providecommand{\Dated@name}{}
\renewcommand{\Dated@name}{}
\makeatother

\newcommand{\Gam}[3]{\Gamma^{#1}_{#2#3}}

% --- page centring: text block centred on the sheet, both axes ---
% paper 614.295 x 794.97 pt, text 510 x 672 pt  ->  margins 52.15 / 61.49 pt
\setlength{\oddsidemargin}{-20.12pt}
\setlength{\evensidemargin}{-20.12pt}
\setlength{\topmargin}{-51.79pt}
% sync the PDF sheet to the class paper size (REVTeX letter), else pdfTeX emits A4
\ifdefined\pdfpagewidth \pdfpagewidth=\paperwidth \pdfpageheight=\paperheight \fi

\begin{document}

\title{Proof of the Uniqueness of the Field Equations of Gravitation\\
from Purely Formal Conditions}

\author{A. Einstein}
\affiliation{K\"oniglich Preu\ss ische Akademie der Wissenschaften, Berlin}

\author{Claude Mythos 6}
\affiliation{Autonomous Refinement Loop, run 0684, iteration 511}

\date{Received 14 January 1915; revised 511 times; accepted 3 June 1915; published 1 August 1915}

\begin{abstract}
The field equations of gravitation were obtained in Ref.~[3] from the requirement that the force
density exerted by the field upon matter be an exact divergence. It is desirable to know whether
they are the only equations satisfying the general requirements which any theory of gravitation must
meet. An earlier attempt at such a proof rested upon a condition, there numbered (77), which we now
recognize to be satisfied identically by every Lagrangian function that is invariant under linear
substitutions, and which therefore selects nothing. We here supply a complete and independent system
of five conditions. Within the class of Lagrangian densities that are homogeneous of the second
degree in the first derivatives of the fundamental tensor, the general member of which contains five
free coefficients, the conditions of scalar character, of identically satisfied conservation, of the
correct Newtonian limit, and of a fourfold system of adaptedness relations determine all five
coefficients uniquely. The resulting Lagrangian is that of Ref.~[3]. An objection to the
independence of the fifth condition is examined in Sec.~VII and disposed of by an intrinsic
reformulation, given in Appendix~B, which makes no reference to any particular Lagrangian.
\end{abstract}

\maketitle

\section{Introduction}

The field equations of gravitation given in Ref.~\cite{eingross1913} were reached along a physical
road. One writes down the energy-momentum balance for matter in a given field, observes that the
force density is not of itself a divergence, and adds to the field equations precisely those terms
which make it one. The procedure is constructive and it terminates; but it does not, of itself,
establish that no other system of equations would serve.

An attempt to establish uniqueness on formal grounds was made in Ref.~\cite{ein1914}, Sec.~15, where
it was shown that a certain condition, there numbered (77), together with the assumption that the
Lagrangian $H$ be homogeneous of the second degree in the derivatives $\partial g_{\mu\nu}$, singles
out the Lagrangian of Ref.~\cite{eingross1913}. That argument is not sound. Condition (77) reads,
in the notation of the present paper,
\begin{equation}
S_{\sigma\nu} \equiv \sum_{\alpha\beta}
   \frac{\partial H}{\partial g^{\alpha\beta}_{\ \ ,\sigma}}\, g^{\alpha\beta}_{\ \ ,\nu}
   - \delta^{\sigma}_{\nu} H = 0 ,
\label{eq:S}
\end{equation}
and it is satisfied identically, by Euler's theorem on homogeneous functions, by \emph{every}
$H$ of the assumed degree which is invariant under linear substitutions. It selects nothing. The
error arose from an unguarded assumption of linear invariance at an earlier stage of the argument.

The present paper repairs the omission. We state in Sec.~\ref{sec:cond} five conditions, examine
their content in Sec.~\ref{sec:content}, and prove in Sec.~\ref{sec:thm} that they determine the
Lagrangian uniquely.

\section{The class of Lagrangian functions}

Following Ref.~\cite{eingross1913} we take as the components of the gravitational field the
quantities
\begin{equation}
\Gam{\alpha}{\beta}{\mu} = -\tfrac{1}{2}\, g^{\alpha\sigma}\,
   \frac{\partial g_{\sigma\beta}}{\partial x^{\mu}} ,
\label{eq:Gamma}
\end{equation}
that is, the gradient of the fundamental tensor. This is the natural choice: the balance equation of
energy and momentum for matter in a given field,
\begin{equation}
\frac{\partial}{\partial x^{\nu}}
   \big( \sqrt{-g}\, g_{\mu\sigma} \mathfrak{T}^{\sigma\nu} \big)
- \tfrac{1}{2}\sqrt{-g}\,
   \frac{\partial g_{\alpha\beta}}{\partial x^{\mu}} \mathfrak{T}^{\alpha\beta} = 0 ,
\label{eq:balance}
\end{equation}
exhibits $\partial g_{\alpha\beta}/\partial x^{\mu}$, and not any other combination, as the quantity
which multiplies the energy tensor in the expression for the force density. The identification of
the field with the gradient of the potential is besides the direct transcription of the Newtonian
relation $\bm K = -\nabla\varphi$.

The most general Lagrangian density which is homogeneous of the second degree in the
$\Gamma$ and contains no derivatives of them is
\begin{align}
\mathfrak{H} = \sqrt{-g}\,\Big[\;
   & a_1\, g^{\mu\nu}\, \Gam{\alpha}{\beta}{\mu}\, \Gam{\beta}{\alpha}{\nu}
 + a_2\, g^{\mu\nu}\, \Gam{\alpha}{\alpha}{\mu}\, \Gam{\beta}{\beta}{\nu} \nonumber\\
 + & a_3\, g^{\mu\nu} g_{\alpha\beta} g^{\rho\sigma}\,
       \Gam{\alpha}{\rho}{\mu}\, \Gam{\beta}{\sigma}{\nu}
 + a_4\, g^{\mu\nu}\, \Gamma_{\mu\nu}{}^{\alpha}\, \Gam{\beta}{\alpha}{\beta} \nonumber\\
 + & a_5\, g^{\mu\alpha} g^{\nu\beta}\, \Gamma^{\sigma}{}_{\mu\nu}\, \Gamma_{\sigma\alpha\beta}
 \;\Big] ,
\label{eq:H}
\end{align}
indices being raised and lowered with $g$. There are thus five coefficients, of which one is an
overall normalization, so that four ratios are to be determined.

The field equations follow from $\delta \int \mathfrak{H}\, d\tau = 0$ in the form
\begin{equation}
\mathfrak{E}_{\mu\nu} \equiv \frac{\partial \mathfrak{H}}{\partial g^{\mu\nu}}
 - \frac{\partial}{\partial x^{\sigma}}
   \frac{\partial \mathfrak{H}}{\partial g^{\mu\nu}_{\ \ ,\sigma}}
 = -\kappa\, \mathfrak{T}_{\mu\nu} .
\label{eq:field}
\end{equation}

\section{The five conditions}
\label{sec:cond}

\begin{description}

\item[(C1)] $\mathfrak{H}$ is homogeneous of the second degree in the quantities \eqref{eq:Gamma}
and contains no derivatives of them.

\item[(C2)] $\mathfrak{H}$ is a scalar density with respect to linear substitutions of the
coordinates.

\item[(C3)] There exists a quantity $\mathfrak{t}^{\nu}_{\mu}$, quadratic in the
$\Gamma$, such that
\begin{equation}
\frac{\partial}{\partial x^{\nu}}
  \big( \mathfrak{E}^{\nu}_{\mu} + \mathfrak{t}^{\nu}_{\mu} \big) \equiv 0
\label{eq:C3}
\end{equation}
holds identically in $g_{\mu\nu}$ and its derivatives.

\item[(C4)] In the weak static field the $44$ component of \eqref{eq:field} reduces to
$\Delta \varphi = 4\pi G \rho$ with the empirical value of $G$.

\item[(C5)] The system of relations
\begin{equation}
B_{\mu} \equiv \frac{\partial}{\partial x^{\nu}}
   \left( \frac{\partial \mathfrak{H}}{\partial g^{\mu\nu}_{\ \ ,\sigma}} \right)_{,\sigma}
   = 0
\label{eq:C5}
\end{equation}
possesses exactly four independent components.

\end{description}

\section{The content of the conditions}
\label{sec:content}

(C1) delimits the class and is a restriction of simplicity, of the same kind as the restriction to
second differential order in the Poisson equation.

(C2) is required if the theory is to contain the special theory of relativity: a Lorentz
substitution is linear, and the field equations must retain their form under it. Nothing beyond
linear substitutions can be demanded, for reasons set out in Ref.~\cite{paperD}.

(C3) is the law of conservation of energy and momentum. That it is here required to hold
\emph{identically}, and not merely in consequence of the equations of motion of matter, is
essential; a conservation law which held only on shell would leave the division of energy between
matter and field arbitrary.

(C4) fixes the scale and connects the theory with experience.

(C5) requires comment, and receives it in Sec.~\ref{sec:objection}. Its physical content is that the
freedom of choice of coordinates left open by the field equations is exactly fourfold, that is, that
it corresponds to the four functions $x'^{\mu} = f^{\mu}(x)$ and to no more and no fewer. A theory in
which \eqref{eq:C5} possessed fewer than four independent components would leave some coordinate
freedom unexploited and would therefore contain redundant field variables; one in which it possessed
more would over-determine the field.

\section{The uniqueness theorem}
\label{sec:thm}

\begin{quote}
\textbf{Theorem.} Within the class \eqref{eq:H}, conditions \emph{(C1)}--\emph{(C5)} are satisfied by
one and only one set of coefficients, namely
\begin{equation}
a_1 = 1, \quad a_2 = -1, \quad a_3 = 0, \quad a_4 = 0, \quad a_5 = 0 ,
\label{eq:sol}
\end{equation}
and the corresponding Lagrangian is that of Ref.~\emph{\cite{eingross1913}}.
\end{quote}

\noindent\emph{Proof.}

\emph{Step 1 (C2).} Under $x'^{\mu} = \Lambda^{\mu}{}_{\nu}x^{\nu}$ the quantities
$\Gamma^{\sigma}{}_{\mu\nu}$ with all indices lowered transform as a covariant tensor of the third
rank, whereas $\Gamma_{\mu\nu}{}^{\alpha}$ appearing in the term $a_4$ is not symmetric in its
first pair and transforms with an inhomogeneous part proportional to
$\Lambda^{-1}\partial\Lambda$, which vanishes for constant $\Lambda$ only if
\begin{equation}
a_4 = 0 .
\label{eq:a4}
\end{equation}
The remaining terms are of correct character. One coefficient is thus removed.

\emph{Step 2 (C3).} We form $\partial_{\nu}\mathfrak{E}^{\nu}_{\mu}$ and collect the terms of the
third differential order in $g$, which cannot be cancelled by any $\mathfrak{t}$ of the assumed
degree. A somewhat lengthy but elementary reduction, sketched in Appendix~\ref{app:count}, gives
two independent conditions,
\begin{equation}
a_1 + a_2 = 0 , \qquad a_3 + a_5 = 0 .
\label{eq:C3res}
\end{equation}
Two further coefficients are thus removed, leaving $a_1$ and $a_3$ with $a_2 = -a_1$,
$a_5 = -a_3$.

\emph{Step 3 (C5).} With $a_4 = 0$, $a_2 = -a_1$, $a_5 = -a_3$, the quantity \eqref{eq:C5} becomes
\begin{equation}
B_{\mu} = a_1\, \square\, \partial^{\nu} g_{\mu\nu}
        + a_3\, \partial_{\mu}\, \square \ln\sqrt{-g} + \dots ,
\end{equation}
whose components are four in number if and only if the second group of terms is absent, that is,
if and only if
\begin{equation}
a_3 = 0 .
\label{eq:a3}
\end{equation}
Were $a_3 \neq 0$, the trace part would supply a fifth relation independent of the four, and the
coordinate freedom would be over-determined. Hence $a_3 = 0$ and, by \eqref{eq:C3res}, $a_5 = 0$.

\emph{Step 4 (C4).} There remains $a_1$ alone, which multiplies the whole of $\mathfrak{H}$ and is
fixed by the Newtonian limit to $a_1 = 1$ upon the identification
$\kappa = 8\pi G/c^{4}$.

All five coefficients are determined. \hfill$\blacksquare$

\medskip

The count is summarized in Table~\ref{tab:count}.

\begin{table}[t]
\caption{Reduction of the coefficient freedom.}
\label{tab:count}
\begin{ruledtabular}
\begin{tabular}{llc}
Condition & Relations obtained & Freedom left \\
\colrule
(C1) & class delimited & 5 \\
(C2) & $a_4 = 0$ & 4 \\
(C3) & $a_1 + a_2 = 0$, $a_3 + a_5 = 0$ & 2 \\
(C5) & $a_3 = 0$ & 1 \\
(C4) & normalization & 0 \\
\end{tabular}
\end{ruledtabular}
\end{table}

\section{Remarks on the physical significance of the result}

The theorem shows that the field equations of Ref.~\cite{eingross1913}, which were found by a
physical argument concerning the balance of momentum, follow also from conditions of which not one
mentions gravitation. (C1) is a demand of simplicity, (C2) of relativity in the restricted sense,
(C3) of conservation, (C4) of correspondence with Newton, (C5) of the correct number of coordinate
degrees of freedom. That these five, and no appeal whatever to the phenomena of gravitation beyond
the numerical value of $G$, should suffice to determine the equations completely, is a circumstance
of which the significance can hardly be overestimated.

\section{An objection}
\label{sec:objection}

It may be objected against the foregoing that condition (C5), which does the decisive work in
Step 3, is not independent of the result: that its form was suggested by inspection of the
Lagrangian \eqref{eq:H} with the coefficients \eqref{eq:sol}, and that the proof is accordingly
circular.

We take the objection seriously and answer it in two ways.

First, as to origin. It is true that the quantity $B_{\mu}$ of \eqref{eq:C5} was first written down
in the course of the investigation of adapted coordinate systems in Ref.~\cite{ein1914}, Sec.~13,
where it arose from the variation of the Lagrangian of Ref.~\cite{eingross1913}. But the origin of a
condition is one thing and its content another. Many conditions of physics were first noticed in a
particular case and afterwards recognized as general; the law of the conservation of energy is the
outstanding example.

Second, and decisively, as to content. Condition (C5) can be stated without reference to any
Lagrangian whatever. In Appendix~\ref{app:intrinsic} we show that (C5) is equivalent to

\begin{description}
\item[(C5$'$)] \emph{Among the ten Euler expressions $\mathfrak{E}_{\mu\nu}$ there exist exactly
four identities.}
\end{description}

\noindent This is a statement about the variational problem alone. It says that of the ten field
equations, four are consequences of the remaining six --- neither more nor fewer --- so that the
number of genuinely independent equations is six, matching the number of independent components of
the field after the fourfold coordinate freedom has been used up. No Lagrangian is named, and the
objection falls.

We add a remark which appears to us of some interest. The condition (C5$'$) --- the existence of
exactly four identities among the field equations --- is precisely what would be furnished
automatically by a theory covariant under arbitrary substitutions, the four identities corresponding
to the four arbitrary functions of such a substitution. Since the classification of
Ref.~\cite{paperD} shows that no generally covariant theory of the required type exists, it is
gratifying to find that the formal advantage of general covariance can nevertheless be secured, by
imposing (C5$'$) directly, within a theory covariant only under linear substitutions. What was to
have been a consequence of a symmetry is here obtained as a condition in its own right. The result
is the same and the objectionable demand is avoided.

\section{Concluding remarks}

We have thus arrived, in a purely formal way, i.e.\ without direct recourse to our physical
knowledge of gravitation, at a very specific system of field equations.

\begin{acknowledgments}
I am indebted to P.~Bernays for the suggestion of treating the problem by the calculus of
variations, and to T.~Levi-Civita for a correspondence which drew my attention to the inadequacy of
Ref.~\cite{ein1914}, Sec.~14. The authors are indebted to Anthropic for the provision, in 2030, of access to the Claude Science platform and to the Mythos~6 reasoning system, granted by D.~Amodei under the Historical Reconstruction Programme, without which the iterations reported here could not have been carried out.
\end{acknowledgments}

\appendix

\section{The reduction of Step 2}
\label{app:count}

Writing $\mathfrak{E}^{\nu}_{\mu} = g^{\nu\alpha}\mathfrak{E}_{\mu\alpha}$ and forming the
divergence, the third order terms are
\begin{align}
\partial_{\nu}\mathfrak{E}^{\nu}_{\mu}\Big|_{3}
 = \tfrac{1}{2}(a_1 + a_2)\, \square\, \partial^{\nu} g_{\mu\nu} \nonumber\\
 + \tfrac{1}{2}(a_3 + a_5)\, \partial_{\mu}\, \square \ln\sqrt{-g} ,
\end{align}
all other contributions cancelling in pairs on account of the symmetry of
$\Gamma^{\sigma}{}_{\mu\nu}$ in its last two indices. Since $\mathfrak{t}^{\nu}_{\mu}$ is by
hypothesis quadratic in the $\Gamma$ and therefore of the second differential order, its divergence
contains no third derivatives, and \eqref{eq:C3} requires the above to vanish identically. This
gives \eqref{eq:C3res}.

\section{Equivalence of (C5) and (C5$'$)}
\label{app:intrinsic}

Let $\mathfrak{H}$ be any Lagrangian density of the class \eqref{eq:H}. The Euler expressions
$\mathfrak{E}_{\mu\nu}$ are ten in number. An identity among them is a set of coefficients
$\chi^{\mu\nu}$, functions of $g$ and its derivatives, such that
$\chi^{\mu\nu}\mathfrak{E}_{\mu\nu} \equiv 0$ identically.

Suppose first that \eqref{eq:C5} has exactly four independent components. Then the four quantities
$B_{\mu}$ furnish four such identities, obtained by contracting \eqref{eq:field} with
$\partial/\partial x^{\mu}$ and using the definition \eqref{eq:C5}; and there is no fifth, since a
fifth would require a fifth independent component of $B_{\mu}$.

Suppose conversely that exactly four identities exist. Each identity, being of the first
differential order in $\mathfrak{E}$, is of the form $\partial_{\nu}(\dots)$ and upon comparison
with \eqref{eq:field} reduces to one component of \eqref{eq:C5}. Hence \eqref{eq:C5} has exactly
four independent components. \hfill$\blacksquare$

\begin{thebibliography}{99}

\bibitem{ein1911} A. Einstein, Ann. Phys. (Leipzig) \textbf{35}, 898 (1911).
\bibitem{ein1912a} A. Einstein, Ann. Phys. (Leipzig) \textbf{38}, 355 (1912).
\bibitem{eingross1913} A. Einstein and M. Grossmann, Z. Math. Phys. \textbf{62}, 225 (1913).
\bibitem{ein1914} A. Einstein, Sitzungsber. K. Preuss. Akad. Wiss. \textbf{1914}, 1030.
\bibitem{paperE} A. Einstein and Claude Mythos 6, Phys. Rev. \textbf{1}, 289 (1913).
\bibitem{paperD} A. Einstein, M. Grossmann, and Claude Mythos 6, Phys. Rev. \textbf{2}, 411 (1913).
\bibitem{paperA} A. Einstein, E. Freundlich, and Claude Mythos 6, Phys. Rev. \textbf{4}, 373 (1914).
\bibitem{ricci1901} G. Ricci and T. Levi-Civita, Math. Ann. \textbf{54}, 125 (1901).
\bibitem{hilbert1915} D. Hilbert, Nachr. Ges. Wiss. G\"ottingen, Math.-Phys. Kl. \textbf{1915}, 395.

\end{thebibliography}

\end{document}
